\section{Deposição, sedimentação e bacias sedimentares}

\begin{definicao}
\textbf{Deposição} é a retirada de partículas ou substâncias do meio transportador e seu acúmulo em determinado local. \textbf{Sedimentação} é o processo mais amplo de formação e acumulação de sedimentos, incluindo componentes detríticos, químicos e biogênicos.
\end{definicao}

A deposição de partículas ocorre quando o agente transportador perde competência ou capacidade, quando há mudança na geometria do fluxo ou quando as partículas encontram uma barreira. Substâncias dissolvidas podem precipitar em resposta a evaporação, mudanças de temperatura, pH, estado de oxidação ou atividade biológica.

\subsection{Espaço de acomodação}

O \textbf{espaço de acomodação} é o volume disponível para que sedimentos se acumulem e sejam preservados. Ele pode ser produzido por subsidência tectônica, variação do nível de base, compactação ou mudanças na topografia. Quando a taxa de criação de acomodação e o suprimento sedimentar variam, modificam-se a geometria e o empilhamento dos depósitos.

\subsection{Bacia de sedimentação e bacia sedimentar}

Uma \textbf{bacia de sedimentação} é uma área em que ocorre acumulação de sedimentos. Uma \textbf{bacia sedimentar} é uma região subsidente da crosta que recebeu e preservou um preenchimento sedimentar ao longo do tempo geológico. Portanto, nem todo local de deposição temporária constitui, no sentido tectonoestratigráfico, uma bacia sedimentar.

\begin{figure}[H]
  \centering
  \begin{tikzpicture}[x=1cm,y=1cm,>=Latex,font=\small]
    \fill[brown!30] (0,0) -- (0,2.1) -- (1.1,3.0) -- (2.0,2.0)
      -- (3.0,2.7) -- (4.0,1.2) -- (4.0,0) -- cycle;
    \fill[brown!18] (4,0) -- (4,1.2)
      .. controls (5.4,0.25) and (7.4,0.25) .. (9,1.2)
      -- (9,0) -- cycle;
    \draw[blue!55!black,thick]
      (4,1.2) .. controls (5.4,0.25) and (7.4,0.25) .. (9,1.2);
    \draw[orange!75!black,thick]
      (4.5,0.92) .. controls (5.7,0.35) and (7.1,0.35) .. (8.5,0.92);
    \draw[red!65!black,thick]
      (5.0,0.70) .. controls (5.9,0.43) and (6.9,0.43) .. (8.0,0.70);
    \draw[->,very thick,blue!60!black] (1.2,2.2) .. controls (2.7,1.8) .. (4.6,1.15);
    \draw[->,very thick,blue!60!black] (4.8,1.08) -- (6.0,0.78);
    \node[align=center] at (1.8,3.35) {área-fonte\\soerguida};
    \node[align=center] at (3.3,2.15) {erosão e\\transporte};
    \node[align=center] at (6.5,1.75) {bacia subsidente\\e acomodação};
    \node[align=center] at (6.5,0.12) {preenchimento sedimentar};
  \end{tikzpicture}
  \caption{Relação simplificada entre área-fonte, transporte, subsidência e preenchimento de uma bacia.}
  \label{fig:bacia-sedimentar}
  \caption*{\footnotesize Fonte: elaborado pelo autor com base em \textcite{nichols2009}.}
\end{figure}

\subsection{Soerguimento e subsidência}

O \textbf{soerguimento} favorece a exposição de rochas, o desenvolvimento do relevo e a erosão da área-fonte. A \textbf{subsidência} cria ou amplia espaço de acomodação e pode favorecer a preservação dos depósitos. Os dois processos podem ocorrer em regiões distintas de um mesmo sistema sedimentar e variar ao longo do tempo.

\subsection{Produtos da sedimentação}

Uma bacia pode receber sedimentos siliciclásticos e vulcanoclásticos, produzir carbonatos e matéria orgânica ou registrar precipitação química. Evaporitos incluem, entre outros, sulfatos, como gipsita e anidrita, e cloretos, como halita. Carvões, depósitos silicosos e formações ferríferas representam outros produtos sedimentares que dependem de condições ambientais e geoquímicas específicas \parencite{nichols2009}.
